\section{Stabilizer and Quasiprobability Decompositions}
\label{app:stabilizer}

This appendix relates common magic measures to the primary local cost \(\xi(A)\). In the theorem, \(\xi(A)\) means a certified free-tensor quasiprobability \(\ell_1\) cost. Other measures are used only insofar as they provide decompositions or upper bounds.

\subsection{Stabilizer rank and extent}

For a pure \(n\)-qubit state \(\ket{\psi}\), the stabilizer rank \(\chi(\psi)\) is the minimum number of stabilizer states needed in an exact linear decomposition
\[
\ket{\psi}=\sum_{j=1}^{\chi} c_j\ket{\phi_j}.
\]
Approximate stabilizer rank permits approximation in norm or fidelity and is used in simulation algorithms where small state error induces controlled output error. Stabilizer extent is the squared minimum \(\ell_1\)-norm of stabilizer decompositions in common conventions:
\[
\xi_{\mathrm{ext}}(\psi)
=\min_{\psi=\sum_j c_j\phi_j}\left(\sum_j |c_j|\right)^2.
\]
Because conventions differ by squares, the main text uses a generic \(\xi(A)\) and places squares explicitly in variance statements.

\subsection{Robustness of magic}

Robustness of magic expresses a state or operation as a signed combination of free stabilizer objects and minimizes the \(\ell_1\) weight. It has an operational role in quasiprobability simulation and magic-state resource theory. For mixed states,
\[
\rho=\sum_j q_j\sigma_j,\qquad \sigma_j\in\Sstab,
\]
with real coefficients \(q_j\), the relevant cost is typically \(\sum_j |q_j|\), minimized over stabilizer decompositions. Channel robustness and related quantities can be defined analogously through free channels or Choi states, with care about normalization.

\subsection{Dyadic negativity and quasiprobability norm}

Quasiprobability simulation chooses a representation in which free operations are efficiently sampleable and resource operations have signed representations. The cost is an \(\ell_1\) norm of the signed representation. Dyadic negativity is one such measure adapted to stabilizer operations over qubits. The SLMS theorem is compatible with any such measure satisfying:
\begin{enumerate}[leftmargin=1.2cm]
\item exact reconstruction of the resource tensor as a signed combination of free tensors;
\item efficient sampling from absolute weights;
\item a second-moment bound controlled by the square of the \(\ell_1\) norm;
\item closure of sampled free tensors under the message-passing contraction.
\end{enumerate}

\subsection{Worked examples for \texorpdfstring{\(\kappa_b\)}{kappa}}
\label{app:kappa-example}

The factor \(\kappa_b\) is a certified local expansion factor for the coefficient map induced by contracting a bag. Fix a finite free dictionary \(\{F_i\}_{i=1}^m\) for the incoming representatives and \(\{G_j\}_{j=1}^{m'}\) for the outgoing representatives. If the local contraction sends
\[
F_i\longmapsto \sum_j (L_b)_{j i}G_j,
\]
then a coefficient vector \(q\) is mapped to \(L_b q\). A sufficient bound on the outgoing \(\ell_1\) cost is
\[
\norm{L_b q}_1\le \norm{L_b}_{1\to 1}\norm{q}_1,
\qquad
\kappa_b=\norm{L_b}_{1\to1}
=\max_i\sum_j |(L_b)_{j i}|.
\]
The message recurrence is therefore
\[
\Gamma_b=\kappa_b\xi_b\prod_{c\in\Act(b)}\Gamma_c .
\]

For a stochastic classical marginalization map, let \(G_0,G_1\) be the two outgoing classical basis messages and suppose the local map is
\[
L_{\mathrm{stoch}}
=
\begin{pmatrix}
1 & 0 \\
0 & 1
\end{pmatrix}
\quad\text{or more generally}\quad
L_{j i}\ge 0,\quad \sum_j L_{j i}\le 1 .
\]
Then \(\norm{L_{\mathrm{stoch}}}_{1\to1}\le 1\), so \(\kappa_b=1\). This is the case for normalized classical postprocessing, erasure, and marginalization maps used in the leaf-local examples.

As a nontrivial expansion example, consider changing from a one-dimensional signed representative \(F_+\) to a two-element outgoing dictionary by
\[
F_+\longmapsto G_0-G_1 .
\]
The coefficient matrix is a single column \(L=(1,-1)^T\), hence
\[
\kappa_b=\norm{L}_{1\to1}=|1|+|-1|=2 .
\]
A bag using this map doubles the certified coefficient norm, and the factor is explicitly charged in \(\Gamma_b\). Thus \(\kappa_b\) is not a hidden assumption: it is a local certificate that must be computed, upper-bounded, or conservatively absorbed by marking a larger region active.

\subsection{Operational measure used in the theorem}

The primary operational measure is the generic free-tensor quasiprobability \(\ell_1\) norm. Stabilizer extent, robustness, and dyadic negativity are treated as instantiations or upper bounds, not as separate theorem-level primitives. Comparing the resulting \(\Lammsg\) certificates under several decomposition families is a separate empirical and algorithmic task.
